%=============================================================================
% MEASUREMENT PROBLEM — LITERATURE CONTEXT (Section 17.7 prefix)
% TKT-2026-05-12-manuscript-elite-revision Step 2/4
%
% Provides explicit literature embedding for the Born-rule research
% programme, as required by the elite-review standard.
% Addresses critique: UIDT must be positioned relative to the 2024-26
% landscape of interpretations before presenting its own candidate.
%=============================================================================
\subsection{The Measurement Problem: Current Landscape (2024--2026)}
\label{sec:mp_context}

The quantum measurement problem---the question of why and how definite
outcomes arise from unitary evolution, and why Born-rule probabilities
take the form $p = |\psi|^2$---remains one of the most actively debated
foundational issues in physics~\cite{Zurek2003,Schlosshauer2004}.
Recent systematic reviews confirm that no interpretation has yet achieved
consensus: decoherence-based approaches~\cite{Schlosshauer2004},
relational quantum mechanics~\cite{Rovelli1996}, QBism~\cite{Fuchs2014},
many-worlds formulations~\cite{Everett1957}, and spontaneous-collapse
theories~\cite{GRW1986} each resolve certain aspects while introducing
fresh conceptual or empirical challenges~\cite{MeasurementReview2025}.

A critical diagnostic from the recent literature~\cite{MeasurementReview2025}
identifies three unresolved sub-problems that any satisfactory account
must address:
\begin{enumerate}[leftmargin=1.8em,itemsep=2pt]
  \item the \emph{definite-outcome problem}: why does measurement produce
        one result rather than a superposition;
  \item the \emph{Born-rule derivation problem}: why probabilities equal
        squared amplitudes rather than some other functional form; and
  \item the \emph{preferred-basis problem}: what selects the basis in
        which outcomes are definite.
\end{enumerate}
Decoherence alone addresses (3) effectively but does not resolve (1) or (2)
without additional assumptions~\cite{Schlosshauer2004}.

\subsubsection*{Position of \textsc{uidt} in this Landscape}

\textsc{uidt} falls within the broad family of \emph{epistemically-grounded,
deterministically-structured} interpretations, with closer structural
affinity to relational~\cite{Rovelli1996} and informational~\cite{Fuchs2014}
accounts than to collapse or many-worlds approaches. The specific
contribution proposed in Section~\ref{sec:born} is a candidate mechanism
for (2)---the Born-rule derivation---via information-grid ergodicity.
Problems (1) and (3) are addressed only partially and are listed as
limitations (L8) in Section~\ref{sec:limitations}.

\begin{tcolorbox}[killswitch]
  Kill-Switch~F7: If no information-grid realisation satisfying the three
  technical conditions T1--T3 (Section~\ref{sec:born}) can be constructed
  such that Born-rule statistics emerge in the thermodynamic limit, the
  \textsc{uidt} measurement proposal is falsified and must be retracted
  from the framework. The phenomenological core (Pillars~I--III) is
  unaffected.
\end{tcolorbox}
